\documentclass[12pt]{article}
\usepackage{amsmath}
\usepackage{mathrsfs}
\usepackage{eucal}
\usepackage{amssymb}
\usepackage[french]{babel}
\usepackage[latin1]{inputenc}
\usepackage[T1]{fontenc}
\usepackage[all]{xy}
%\usepackage{graphicx}
%\usepackage{epsf}

\begin{document}

\hfuzz 8 pt
\parindent = 0 cm

\centerline{ \textsc{ \large Universit� de Sherbrooke}}
\centerline{\textsc{ \normalsize D�partement de math�matiques}}
%\centerline{\textsc{ \normalsize Automne 2005}}

\bigskip

\centerline{\large\bf MAT501 -- Fondements et histoire des
math�matiques}

\bigskip
\centerline{\Large\bf Laboratoire 3}
\bigskip


\begin{description}

\item[1)] Montrer que l'ensemble vide est unique et qu'il est
contenu dans tous les ensembles.


\item[2)] L'axiome de r�gularit� s'�nonce ainsi

\begin{verse}
Si $x$ est un ensemble non vide, alors il existe un ensemble $y
\in x$ tel que $t \in x$ implique $t \notin y$.
\end{verse}

Autrement dit, $x \cap y = \varnothing$. Noter que cet axiome ne
dit pas que tout $y \in x$ est disjoint de $x$, mais qu'il y en a
au moins un.

    \begin{description}

    \item[a)] Montrer que l'axiome de fondation implique l'axiome
    de r�gularit� ;

    \item[b)] Montrer que l'axiome de r�gularit� implique que $x
    \notin x$ pour tout ensemble $x$ ;

    \item[c)] En conclure que cet axiome implique qu'il n'existe
    pas un ensemble de tous les ensembles ;

    \item[d)] En d�duire que la collection $R$ des ensembles $x$
    avec $x \notin x$ n'est pas un ensemble (c'�tait le paradoxe de
    Russell).

    \end{description}


\item[3)] Montrer que l'axiome de fondation implique qu'il est
impossible d'avoir des ensembles $x$ et $y$ avec $y \in x \in y$.
Peut-on g�n�raliser?


\item[4)]

    \begin{description}

    \item[a)] Montrer que l'ensemble vide n'est le successeur d'aucun
ensemble.

    \item[b)] Montrer que des ensembles diff�rents ont des
    successeurs diff�rents.

    \end{description}


\item[5)] Soit $A$ un ensemble. Montrer que la collection des
ensembles �quipotents � $A$ ne peut former un ensemble.
[\emph{Indication :} Proc�der par l'absurde en supposant que c'est
un ensemble $B$ et consid�rer $\mathcal{P}(B)$ ainsi que
l'ensemble $U$ des $A_C$ o� $A_C = \{ (x,C) \mid x \in A \}$ pour
$C \in \mathcal{P}(B)$.]

\end{description}

\end{document}
